\section{What changes in v9}
\label{sec:short:changes}

Six changes, each with a date. Every one of them is planned rather than shipped unless the sentence
says otherwise, and the reader should hold that distinction throughout: this section describes a
network being rebuilt, and the full paper marks every element of it that is still undecided.

\subsection{Operators are paid by demand, not by issuance}

This is the change the others serve. Today an operator's income is new \flux{} minted every block;
the applications running on the network contribute nothing to it. Under \emph{Progressive Node
Rewards} that reverses: consensus takes \SI{20}{\percent} of every application payment on arrival
and pays it into a pool, and the pool drips into operator rewards over a rolling window of
\num{86400} blocks --- thirty days. The other \SI{80}{\percent} goes to the Foundation address that
already receives application payments today.

The split is applied by consensus at the address, not by any software a tenant or a node runs. The
pool balance is committed in every block and recomputed by every validator, so it cannot be
misreported. Activation is at block \textbf{\num{3787502}}, expected \textbf{1 July 2027}. That
height is not arbitrary: it is the block at which the parallel-asset mining allocation is exhausted,
computable today from the emission schedule, though no consensus constant encodes it yet; the
calendar date follows from the height at the thirty-second block target.

A reader should understand what this does and does not change. It changes \emph{where operator
income comes from} --- from issuance toward the workloads themselves. It does not, by design, pay a
node more for hosting more: every eligible node in a tier draws the same share of the pool
regardless of what it runs. The team chose that deliberately, because paying by placement would
turn the first-to-claim rule of \Cref{sec:short:how} into a race with money on it. The consequence
is stated in \Cref{sec:short:limits}.

\subsection{Ten parallel assets become two, on one redesigned bridge}

\Flux{} exists as a token on ten other chains. That footprint is being consolidated. Ethereum and
BNB Chain survive: balances there are carried across to new contracts by snapshot, and holders do
nothing. The other eight --- Base, Polygon, Avalanche, Ergo, Algorand, Solana, Tron and Kadena ---
are retired in stages: from March 2027 (block \num{3450000}) they become one-directional, with balances redeemable back to
the main chain and mining claims still collectable, but nothing swappable into them; the allocation
is cut in July, at the block where
it is exhausted and PNR activates; and in August (block \num{3880000}) the chains shut down, after which
unclaimed balances go to the Foundation. Ethereum and BNB Chain move to their new contracts earlier. What the height does
not settle is the balance already earned: roughly a quarter of what operators have accrued on these
chains has never been claimed, and the forfeiture rule is what closes that.

The numbers were measured rather than assumed, and they are reassuring in a way the raw supply
figures are not. Across all ten chains, \num{223027781.21}~\flux{} is actually held outside the
issuer's own addresses. Of that, \SI{97.21}{\percent} sits on the two surviving chains and needs
no holder action. Only \textbf{\num{6226900.80}~\flux{} --- \SI{2.79}{\percent}} --- is on the
retiring chains and exposed to the window.

The new bridge is mint-and-burn: minting on a parallel chain requires locked backing on the main
chain, so the total can never exceed what exists. It is guarded by a 3-of-4 signing quorum, a
1-of-4 emergency pause, a per-chain cap of \num{250000000}~\flux{}, a rate limit of
\num{10000000}~\flux{} per chain per day, and a 48-hour delay on any upgrade during which a
single guardian can cancel it. The quorum is held by four named people from one organisation,
and the paper says so.

\subsection{A new application specification}

Version 9 of the specification adds what the rest of this section needs: a duration expressed in
seconds rather than blocks, so lifetimes stop moving when block timing changes; a payment memo that
binds a payment to the application it funds, which is what lets PNR attribute revenue; and a
machine-type field that lets a specification ask for accelerated hardware. Enforcement is at block
\textbf{\num{3050000}}, expected in \textbf{late October 2026}, with existing applications migrated
automatically in the same quarter.

\limitnote{The fields are in the schema and not yet in the code that would act on them: as of the
survey, the payment memo, referral code and machine-type fields appear nowhere in the placement
path. The enforcement height is constrained to follow the implementation, not precede it.}

\subsection{The network can remove an application}

A collateral-weighted vote. Any node operator can file a report against an application on chain;
other operators vote by signing a transaction with their collateral key; if votes reach
\SI{25}{\percent} of the network's weight within a day, the next block carries a ban memo and every
node stops the application. The ban is final, the unused prepaid balance is burned, and the deployer
may return only as a new application. There is no appeal.

The team's stated position is that this is a network-level veto over content, not merely over
behaviour. That is a stronger power than any comparable network exercises, chosen knowingly, and its
principal risk is measured rather than hypothetical --- see \Cref{sec:short:limits}.

\subsection{Nodes that can prove what they are}

Two layers. The first is shipped: \ArcaneOS{}, the operating image nodes run, boots through a
measured chain --- Secure Boot, a pinned platform key, a hardware root, encrypted storage --- and
refuses to come up if any link fails. The second is planned: a bonded attestation network in which
nodes stake collateral on claims about what they serve, and lose it if a challenger proves them
wrong. Its parameters are set; its slashing construction, the paper is candid, is unsolved.

\subsection{The shielded pool is retired, and what that costs}

\Flux{}'s chain inherited Zcash's private-transaction machinery. It has been sealed to new entrants
since 2021 and today holds \num{96479.94}~\flux{} --- two hundredths of one percent of supply. The
decision is to retire it: announced now, sunset at block \num{3600000} in April 2027, before PNR
activates, and then remove the code. That makes the chain smaller, removes an inherited trusted
setup, and makes supply exactly auditable. It also means \textbf{\Flux{} will have no chain-level
privacy}: every payment will be publicly linkable, as on Bitcoin. The paper treats payment privacy
as an open problem from that point, not a delivered feature.

\subsection{And one change the fleet's own design forces}

A post-quantum note that is specific to \Flux{}. Because registering a node publishes the
collateral's public key on chain, \num{88514500}~\flux{} --- \SI{20.61}{\percent} of everything
issued --- sits in outputs a sufficiently capable quantum computer could spend immediately. Bitcoin's
comfort, that an unspent output hides its key, does not apply here. The adopted answer uses the one
lever \Flux{} has and Bitcoin does not: node collateral is a registered set that must keep
transacting to be paid, so migration to a post-quantum key can be made a condition of eligibility.
Operators who do not migrate stop earning; nothing is ever frozen.
